\cotefile{156-4}
\scanpage{74}\mspage{73}

\deleted{\uline{Corollaire} 2 \unsure{La}}

\uline{Definition} On dit \deleted{\unsure{alors}} \inserted{que $G$ \uline{raffine} $F$} \deleted{\unsure{dit|si} \unsure{est un} \unsure{raffint|raffinement}} ou est un \uline{raffinement} de $F$ \deleted{. On écrit \unsure{alors}} si $\deleted{\unsure{\in}}\exists$ un $F'$ avec \deleted{$G\ll F$} $G\leq F'\preceq F$. On écrit $G\ll F$

\uline{Corollaire} \deleted{\unsure{2}} la relation $G\ll F$ est une relation d'ordre (impliquée par $G\leq F$ et par $G\preceq F$)
\margin{\uline{Utilise $C6$,} ni \uline{$C5$ ni $C7$}.}

Transitivité : Soit $F\deleted{\unsure{\leq|\preceq}}\ll G\ll H$, donc

\diagram{\begin{tikzcd}[ampersand replacement=\&,row sep=0.1em,column sep=0.3em]
 {} \& \deleted{G} \& {} \& \deleted{H'} \& \deleted{\preceq} \& \deleted{H} \\
 {} \& {} \& {} \& \deleted{\rotatebox{90}{$\leq$}} \& {} \& {} \\
 \deleted{H} \& \deleted{G'} \& \deleted{\preceq} \& \deleted{G} \& {} \& {} \\
 {} \& \deleted{\rotatebox{90}{$\leq$}} \& {} \& {} \& {} \& {} \\
 {} \& \deleted{\unsure{H|G}}\deleted{F} \& {} \& {} \& {} \& {}
 \arrow[from=1-2,to=3-1,dash]\arrow[from=1-4,to=4-2,dash]
 \arrow[from=1-5,to=3-3,dash]\arrow[from=1-6,to=3-5,dash]
\end{tikzcd}}

\diagram{\begin{tikzcd}[ampersand replacement=\&,row sep=0.08em,column sep=0.22em]
 H \& \succeq \& H' \& {} \& {} \\
 \vdots \& {} \& \rotatebox{90}{$\leq$} \& {} \& {} \\
 {} \& \cdots \& G \& \succeq \& G' \\
 {} \& {} \& {} \& {} \& \rotatebox{90}{$\leq$} \\
 {} \& {} \& {} \& {} \& \deleted{\unsure{H}}F
\end{tikzcd}}

on complète par $\unsure{C6|6\text{b}}$

\diagram{\begin{tikzcd}[ampersand replacement=\&,row sep=0.08em,column sep=0.22em]
 H \& \succeq \& H' \& {} \& {} \& {} \\
 \rotatebox{90}{$\leq$} \& \circ \& \rotatebox{90}{$\leq$} \& {} \& {} \& {} \\
 \bar H \& \succeq \& G \& \succeq \& G' \& \quad \raisebox{0.5em}{$(*)$} \\
 \rotatebox{90}{$\leq$} \& \circ \& \rotatebox{90}{$\leq$} \& \circ \& \rotatebox{90}{$\leq$} \& {} \\
 K \& \succeq \& \bar G \& \succeq \& F \& {}
\end{tikzcd}}

et on gagne, on
\[
 F\preceq K\leq H
\]

Supposons maintenant que $F=H$ i.e. $F\ll G$ et $G\ll F$, prouvons $F=\deleted{H}G$. On \deleted{\unsure{tire de}}
\[
 \begin{aligned}
 F\preceq K\leq{} & F & \text{on va en déduire}\\[-2pt]
 &\rotatebox{90}{$=$}\\[-2pt]
 &H & K=F
 \end{aligned}
\]
et comme $\underset{\substack{\rotatebox{90}{$=$}\\F}}{K}\leq\bar H\leq F$ et on \unsure{résultera|resultera|obtiendra}
\[
 \bar H=F\quad \deleted{\text{et de }\overline{\mathrm m}\;
 \unsure{\bar G|\bar H}=\unsure{F|K}\ (\text{comme }F\preceq\bar G\preceq
 \underset{\substack{\rotatebox{90}{$=$}\\F}}{K})\ \bar G=F}
\]
donc $G\preceq F$. De façon symétrique on aura $F\preceq G$, d'où $F=G$. Donc il reste à prouver le

\uline{Lemme} Soient $F,K\in\mathcal F$ avec $F\preceq K\leq F$. Alors $K=F$ (i.e. si $F$ est une subdivision d'une sous-figure $K$, alors $F=K$)

\scanpage{75}\mspage{74}

Cela \unsure{résultera|résulte} du Corollaire qui suit, et de $C\unsure{5|\deleted{6}5}$

\uline{Corollaire} (Propriété de la relation $\leq$ !) Soient $F$ une figure, $K$ une sous-figure \deleted{qui} telle que $K\neq F$. Alors il existe une \deleted{\unsure{str}} figure $G$ (une multistrate, si on veut) qui est \deleted{\unsure{donc|des}} disjointe de $K$, mais non de $F$. En d'autres termes, $\operatorname{supp} K\neq\operatorname{supp} F$.

Ce n'est pas ``formel'' en termes des axiomes $C1 : C4$ sur $(\mathcal F,\leq)$, Supposons p.ex. que $\mathcal F$ admette un plus grand élément, donc il est défini comme l'ens de \uline{toutes} les parties fermées de l'ens ordonné $M$, et ``disjoint'' signifie ici $A\cap B=\phi$ sans plus (la relation de compatibilité $R$ est la relation tjrs vraie). Mais si $F$ est une partie fermée, $K\subset F$ une sous partie fermée \inserted{$\neq F$}, il n'est pas clair qu'on puisse trouver une partie fermée $G$ telle que $G\cap\deleted{F}K=\phi$, $G\cap F\neq\phi$. \deleted{P.ex. $M\unsure{=|-}[0,N]$, $\unsure{F|R}$ avec \unsure{rel\textsuperscript{n}} \unsure{d'ordre|l'ordre} \unsure{(?)|ord} habituel, $F=M$, $K=[0,1]$,} \inserted{Si p.ex. $M$ a un plus petit \unsure{élément|élém\textsuperscript{t}|élt} $0$,} Alors pour \uline{toute} partie fermée non vide, les $G$ tels que $F$ et $G$ disjoints sont les $G$ vides, donc ça ne change pas si on remplace $F$ par $\deleted{\unsure{F|R}}K\in F$, \unsure{tant|\deleted{tant}} que $K$ n'est pas vid\unsure{e}

\scanpage{76}\mspage{75}

Dans ce cas, l'axiome $C6$ (pour une relation \inserted{d'ordre} suppl $\preceq$ sur $\mathcal F=\unsure{\mathrm P|\mathcal P}_{f}(M)$) signifie \unsure{seulement|seulmt} que si $F'\preceq F$, $F'$ est vide ssi $F$ l'est. Prenons p.ex. \deleted{$\unsure{F|K}\preceq$} $\preceq$ \deleted{\unsure{égal|égale}} à la relation d'ordre discrète, c'est OK\unsure{.|,} \deleted{$C7_{\unsure{0|a}}$ est \unsure{tautologique}} $C6,C7$ tautologiques. On voit \unsure{donc|bien} qu'il manque un axiome, pour assurer la validité du corollaire. On devine que ce sera la partie \unsure{--)|a)} du \unsure{Const|Cont} $4$ (p. 35), \unsure{jointe} à Const $7$. Donc pour l'instant \unsure{je} ne \unsure{vais|sais} pas utiliser que le préordre $\ll$ soit un ordre.
\margin{Aussi, cela résulte \unsure{aussitôt|aussi} \unsure{de la} \unsure{modification} de $C4$ ancien (page $47$)}

\deleted{\unsure{Je} songe à la \unsure{propté|propriété} \unsure{suiv\textsuperscript{te}|suiv.|suivante}, qu'il faudrait avoir tot ou tard :}

\deleted{\deleted{\unsure{Soient|Soit} $\unsure{Y,|\forall}X\in M$ \unsure{si|avec|il} $Y\ll X$, et} supposons que \deleted{$Y'\leq Y$, $Y'\in\partial X$, $Y\cap\partial X\unsure{=|-}$} \deleted{$\unsure{\exists|\nexists}\,Y'\leq Y$, \unsure{alors|les}} $Y'\leq Y$, $(Y'\in M)$ tels que \deleted{$\unsure{\phi|\partial|\exists}\,Y'\ll\partial X$}}

\deleted{Soient $F$ figure, $F'$ \unsure{ss-figure|s-figure}, $G\ll F$, $G'$ \inserted{$\leq G$} le \unsure{raffint|raffinement} de $F'$ induit par $G$ cf plus \unsure{bas|pas}). \deleted{\unsure{Supposons|supposons} que} \inserted{Alors} $G'=\unsure{\phi|\varnothing}$. \deleted{Alors $G$ est} \unsure{ssi|ss} $G$ disjoint de $F'$.}

Il faut développer la notion de \uline{\unsure{raffinements|raffinement|raffinants} induit}. J'ai envie de \deleted{\unsure{démontrer|Il en résulte|Il en est}} \inserted{\unsure{prouver|posons}} l'énoncé (quitte à rajouter des axiomes) :

\scanpage{77}\mspage{76}

\deleted{\uline{\unsure{Prop.|Déf.}}} \inserted{\uline{Prop.}} Soit

\diagram{\begin{tikzcd}[ampersand replacement=\&,row sep=0.15em,column sep=0.3em]
 G \& \ll \& F \\
 \deleted{\rotatebox{90}{$\leq$}} \& {} \& \rotatebox{90}{$\leq$} \\
 \deleted{G'\,.} \& {} \& F'\,\unsure{.}
\end{tikzcd}}

alors il existe un plus grande sous figure $F'$ de $G$, qui raffine $F'$ (et qu'on appelle raffinement \unsure{de|en} $F'$ induit \deleted{\unsure{par le|par la}} \unsure{sur la|sur $G$|par $G$} sous figure $F'$ par le raffinement $G$ de $F$, ou la sous-figure du \deleted{$G$ \unsure{induit}} raff $G$ de $F$, induite par la sous-figure $F'$ de $F$). \deleted{\unsure{Elle est}} Les strates \deleted{(c'est $\pm$ tautologique)} sont les strates de $G$ qui raffinent $F'$. \deleted{De plus, $G$ et $F$}

\diagram{\begin{tikzcd}[ampersand replacement=\&,row sep=0.12em,column sep=0.3em]
 G \& \ll \& F \\
 \rotatebox{90}{$\leq$} \& {} \& \rotatebox{90}{$\leq$} \\
 G' \& \ll \& F'
\end{tikzcd}}

De plus, $G'$ est un $\operatorname{Inf}$ de $G$ et $F'$ pour la relation $\underset{\unsure{\cdot}}{\ll}$, en d'autres termes, si \deleted{\unsure{Si|$X$} $G$ est} $K$ est une figure telle que $K\ll G$ et $K\ll F'$, alors $K\ll G'$.

\margin{Dire (\unsure{ou poser|on poser}) que $G'=\phi\Longleftrightarrow G$ disjoint de $F'$ (seul $\Rightarrow$ non évident....)}

\deleted{\uline{\unsure{Dém|Démon}} Soit $H$ une sous-figure \deleted{de $G$ telle que} \inserted{$H$} $H\unsure{\not\ll|\ll} F'$, alors pour \unsure{une|tte} strate $X\leq\deleted{\unsure{F}}H$, on aura $X\ll F'$, donc \unsure{bien|Si} \unsure{on désigne|engendré|engendre} par $G'$ la sous-figure de $G$ \unsure{engendrée|engendré|engendree} par des strates $Y\leq G$ telles que $Y\ll H$, alors \deleted{\unsure{et|$X$}} $H\leq G$ puisque}
\[
 \deleted{\displaystyle H=\mathop{\mathrm{Sup}}\limits_{\substack{X\in\deleted{\unsure{H|G}}\,M\\X\leq H}}X\deleted{\unsure{<}}\leq\mathop{\mathrm{Sup}}\limits_{\substack{X\in M\\X\in G'}}X=G'\unsure{.|,}}
\]
\deleted{D'autre part, $G'\ll$}

\scanpage{78}\mspage{77}

On a par definition de $\ll$
\[
 G\leq\bar F\preceq F
\]
\margin{$G$ sous-figure d'une subd. $\bar F$ de $F$}
et on va appliquer $C6$

\diagram{\begin{tikzcd}[ampersand replacement=\&,row sep=0.07em,column sep=0.25em]
 G \& \leq \& \bar F \& \preceq \& F\\
 \rotatebox{90}{$\leq$} \& {} \& \rotatebox{90}{$\leq$} \& {} \& \rotatebox{90}{$\leq$}\\
 G' \& \leq \& \bar F' \& \preceq \& F'\\
 \rotatebox{90}{$=$} \& {} \& {} \& {} \& {}\\
 G\cap\bar F' \& {} \& {} \& {} \& {}
\end{tikzcd}}

\deleted{\unsure{Soit}} Prouvons que $G'$ est un $\inf$ pour $\ll$

Soit $K$ telle que $K\ll G$, $K\ll F'$ on doit prouver $K\ll G'$. Pour ceci, on prouvera d'abord $K\ll\bar F'$, et on est ramené aux deux cas particuliers où on \unsure{a|-} soit $G\preceq F$, soit $G\leq F$. Donc

\uline{Lemme} 1 Soit

\diagram{\begin{tikzcd}[ampersand replacement=\&,row sep=0.08em,column sep=0.25em]
 G \& \preceq \& F \& \quad\text{(cf } C6\text{)}\\
 \rotatebox{90}{$\leq$} \& {} \& \rotatebox{90}{$\leq$} \& {}\\
 G' \& \preceq \& F' \& {}
\end{tikzcd}}

et soit $K$ \unsure{tel|tels|donc} que $K\ll G$ et $K\ll\deleted{\unsure{G}}F'$, alors $K\ll G'$.

En effet, par \deleted{\unsure{le}} \inserted{def} on aura $K\leq\bar G\preceq\deleted{\unsure{F}}G$

\diagram{\begin{tikzcd}[ampersand replacement=\&,row sep=0.08em,column sep=0.25em]
 K \& \leq \& \bar G \& \preceq \& \deleted{F}G \& \preceq \& F\\
 {} \& {} \& \rotatebox{90}{$\leq$} \& {} \& \rotatebox{90}{$\leq$} \& {} \& \rotatebox{90}{$\leq$}\\
 {} \& {} \& \bar G' \& \preceq \& G' \& \preceq \& F'
\end{tikzcd}}

pour prouver $K\ll G'$, il suffit de prouver $K\leq\overline{G'}$. \deleted{\unsure{Or|donc}} \deleted{Soit $X\in M$ $X\in\unsure{G|C}$} \inserted{ou encore pour} $X\in M$,
\[
 \underset{\inserted{\text{d'où } X\leq\bar G,\ X\ll F'}}{\underline{\rule[-.5ex]{.3pt}{.45ex}\,X\leq K\,\rule[-.5ex]{.3pt}{.45ex}}}\ \Rightarrow X\leq\bar G',
 \quad\text{\unsure{or|ou|un|en}}
\]
Or cela résulte de $C6$, $(b)$ appliqué à

\diagram{\begin{tikzcd}[ampersand replacement=\&,row sep=0.08em,column sep=0.25em]
 \bar G \& \preceq \& F\\
 \rotatebox{90}{$\leq$} \& {} \& \rotatebox{90}{$\leq$}\\
 \bar G' \& \preceq \& F'
\end{tikzcd}}

\scanpage{79}
\mspage{78}

\uline{Lemme 2} Soit
\diagram{\begin{tikzcd}[ampersand replacement=\&,row sep=0.06em,column sep=0.20em]
G \& \leq \& F\\
\rotatebox{90}{$\leq$} \& \& \rotatebox{90}{$\leq$}\\
G' \& \leq \& F'\\
\rotatebox{90}{$=$} \& \& \\
G\cap F' \& \&
\end{tikzcd}}
et soit $K$ tel que $K\ll G$ et $K\ll F'$, alors $K\ll G'$.

On écrit encore $K\leq\bar G\preceq G$, et \ldots

\deleted{\unsure{on}} J'abandonne ! Il vaut mieux se reporter au texte page 51.

\margin{\uuline{NB} $F'\mathrel{\vert\!-\!\vert}G'$ \quad $\text{Sous fig }F-\text{ Sous fig }G$ commute aux $\operatorname{Sup}_{\leq}$ \unsure{quelc}, \deleted{\unsure{puis}} et aux \deleted{\unsure{Sup|Inf}} $\operatorname{Inf}_{\leq}$ quelcon-ques\ldots}

\deleted{Soit $F$ une figure}

\deleted{\uline{\unsure{C8|Co\,8}} Soit $F$ une figure, et $X$ une multistrate telle que $X\ll F$ (i.e $X\leq F'\preceq F$). Alors $\exists\,Y\,\unsure{\in|i}$, \deleted{\unsure{multi}}\unsure{strate} de $F$, telle que $X\deleted{\unsure{\leq}}\ll Y$.}

\uline{Proposition} Soient $F$ une figure, $X\in M$, tels que
\[
 X\ll F.
\]

a) \deleted{Alors, \unsure{je} $\exists\,Y\in$} $\exists\,Y\in M$, $Y\leq F$ telle que
\[
 X\deleted{\unsure{\leq}}\ll Y
\]

b) Parmi les $Y$ satisfaisant \unsure{cette|à cette} condition, il y un plus petit élément, \deleted{\unsure{donné|donc}} \deleted{\unsure{il est|ce}} \inserted{\unsure{J|je|cf|$\hookleftarrow$} \unsure{C avec|conj|avec|car} $C12$.} \deleted{caractérisé par} caractérisé par la propriété que la sous-figure de $X$ induite par $\partial Y$ est $\neq X$.

\uline{Demonstration} a) $F=\mathop{\mathrm{Sup}}\limits_{Y_i\in\bar F}\deleted{\unsure{X_i|Y_i|X|F}}Y_{i\deleted{\unsure{\in F|\in G|\leq|i}}}$\unsure{.|,} donc \deleted{$X$} (par $C6$ $(c)$) $X=\mathop{\mathrm{Sup}}\limits_{i}G_i$ où $G_i$ est \unsure{la|le} \unsure{sous|ss} fig. de $X$ induite par les $\deleted{\unsure{X_i|G_i|F_i}}Y_i$, donc $\exists\,i$ tel que $X\in G_i$, donc $X\ll Y_i$.

\scanpage{80}
\mspage{79}

Parmi les $Y$\inserted{$\in\bar F$} tels que $X\ll Y$, choisissons un élément minimal (cela existe, si on admet que la dim. combinatoire des figures irréd (\unsure{i.e.|$=$.} multistrates) est finie -- on y reviendra sous forme \unsure{d'un|ds un} \unsure{énoncé|environ-t} plus. fort encore \ldots\ \deleted{\unsure{Donc} \unsure{Donc !}}

Dire que $Y$ est minimal signifie que pour $Y'$ \deleted{\unsure{multi|modèle}}strate de $Y$, on n'a pas $X\deleted{\unsure{\leq}}\ll Y'$ i.e. $X\neq F_{Y'}$, sous-figure de $X$ induite par $Y'$. Mais cela signifie aussi que l'image inverse de $\partial Y=\mathop{\mathrm{Sup}}\limits_{\substack{\deleted{\unsure{Y'\leq\partial Y|Y'\leq Y|Y'\in M\leq Y}}\\Y'\in M\\Y'\leq Y}}Y'$ est différent \unsure{de|\deleted{de}} $X$ (appliquant à nouveau $C6$ $(c)$). \deleted{Si on a}

Si on a $Y'\in\bar F$ tel que $X\ll Y'$, prouvons que $\deleted{\unsure{X|\leq}}Y\leq Y'$. \deleted{\unsure{On}} En \deleted{\unsure{appliquant|appliquer} \unsure{le}} effet l'image inverse $F_{Y'}=X$, \unsure{or|on|n}
\[
 \underset{\inserted{\substack{\rotatebox{90}{$=$}\\X}}}{F_Y}\cap
 \underset{\inserted{\substack{\rotatebox{90}{$=$}\\X}}}{F_{Y'}}=F_{Y\cap Y'}
\]
à cause des propriétés formelles du $\operatorname{Inf}$ pour $\ll$ donc $F_{Y\unsure{.|\cdot}\cap Y'}\deleted{\unsure{F_Y\cap F_{Y'}|F_{Y'}|\delta X\cap F_{Y'}}}=X$, i.e. $X$ se factorise par $Y\cap Y'$ donc par un $Z\leq Y\cap Y'\leq Y$ donc $Z=Y$ (\unsure{con.tion|condition} minimum) et par suite $Y=Y'$, on gagne !!

\scanpage{81}
\mspage{80}

\uline{Corollaire} Soient $F,G$ avec $F\ll G$. Pour tout $X\in\bar F$, soit $\varphi(X)\in\bar G$ la plus petite \deleted{multi}strate \inserted{$Y$} de $G$ telle que $\deleted{\varphi(X)\ll}\,\unsure{.}\,X\ll Y$. L'application
\[
 \varphi:\bar F\longrightarrow\bar G
\]
ainsi obtenue est \uline{croissante} (trivial).

Jusqu'à \unsure{maintnt|maintenant}, les axiomes $C5$, $C6$, $C7$ introduits pour $\preceq$ \deleted{\unsure{pour}} sont tautologiques, si on prend pour $\preceq$ l'égalité, \inserted{alors $\ll$ est $\leq$, \unsure{i.e.|scil.} l'inclusion de figures.} Il serait temps d'en venir à des propriétés \deleted{telles que} plus spécifiques !

\uline{14 juin}

\margin{Je vais poser $\widetilde F$ ald $\bar F$ pour l'ens des $X\in M$ tels que $X\leq F$ (pour éviter confusion avec notion d'\unsure{adhérence|adhésion})}

\deleted{\unsure{\uuline{Proposition}|\uline{Proposition}}} \inserted{Corollaire de prop. \unsure{pag.|page} 76} Cf l'\uline{ancien} $C4$ \deleted{\unsure{pa|pag|p}} page 47)

Soient $G'\deleted{\unsure{\leq}}\ll G\deleted{\unsure{\leq}}\preceq F$ alors $\widetilde{G'}\cap\widetilde F\subset\widetilde G$
\inserted{(donc $\widetilde G'\cap\widetilde F=\widetilde{G'}\cap\widetilde G$)}

En d'autres termes, toute \inserted{(multi)}strate \deleted{\unsure{de}} \deleted{\unsure{la} \unsure{figure}} \inserted{$X$ \unsure{(}qui raffine la sous-figure $G$ de $F$ \unsure{et|)}} est \inserted{(multi)}strate de $F$ est \inserted{(multi)}strate de $G$.

\deleted{\deleted{Soit $X$ strate de $G'$ qui est strate de $F$}}

\deleted{\deleted{\unsure{D}} \unsure{Dém.|Dem.} $X$ est une strate d'un raffinement $\bar G$ de $G$, \inserted{lequel est} induit par un raffinement $\deleted{\unsure{G}}\unsure{\bar F|F}$ de $F$ ($C6$)}
\diagram{\begin{tikzcd}[ampersand replacement=\&,row sep=0.05em,column sep=0.16em]
\deleted{G} \& \deleted{\unsure{=|\leq}} \& \deleted{F} \& \deleted{\geq} \& \deleted{X}\\
\deleted{\rotatebox{90}{$\preceq$}} \& \& \deleted{\rotatebox{90}{$\preceq$}} \& \& \deleted{\unsure{;}}\\
\deleted{\bar G} \& \deleted{\leq} \& \deleted{\unsure{.|\cdot}\bar F} \& \& \\
\deleted{\rotatebox{90}{$\leq$}} \& \& \& \& \\
\deleted{X} \& \& \& \&
\end{tikzcd}}
\margin{\unsure{hi hi hi|si si si|hé hé hé|q e.d.|q f.d.} $\longrightarrow$}
\deleted{Supposons $X\leq F$, \deleted{et considérons les \unsure{images|image} \unsure{inverses|inverse}} \inserted{l'image inverse de}}

\scanpage{82}\mspage{81}

\diagram{\begin{tikzcd}[ampersand replacement=\&,column sep=0.6em,row sep=0.15em]
\deleted{X\ll} \& \deleted{G} \& \deleted{\leq} \& \deleted{F}\\
\& \rotatebox{90}{$\deleted{\leq}$} \& \& \rotatebox{90}{$\deleted{\leq}$}\\
\& \deleted{G_X} \& \deleted{\leq} \& \deleted{X}
\end{tikzcd}}

\deleted{Soit $G_X$ la sous figure de $G$ image inverse de \deleted{\unsure{la|les}} sous figure $X$ de $F$. Par la prop. (p.~76) $X$ raffine \unsure{à} : la fois la sous-figure $X$ de $F$, et $G$ raffine $G_X$ \unsure{donc|dans} $X^{\unsure{2|<}}\ll G_X <$}

\deleted{Soit} Supp.\deleted{\unsure{sons|ons|donc}} $X\leq F$, \unsure{i.e.|Id e.|ie.|de.} $X$ ss-fig. de $F$, considérons ses images inverses \inserted{$\bar F_X,\deleted{\unsure{\bar G|G}}G_X,\bar G_X$} dans \deleted{\unsure{les.|les}} $\bar F,G,\bar G$, \inserted{ce sont des sous-figures de $\bar F,G,\bar G$ \unsure{resp.|respt.}} Donc on trouve \deleted{\unsure{(par la prop.|(par le prop.}}
\[
\deleted{X\ll\bar G_X}
\]

\deleted{et on est ramené au cas particulier : $H$ ($=\bar G$) figure, $H'$ ($=\bar G_X$) sous-figure, $X$ multistrate de \deleted{$H$ (\unsure{i.e.|=} $X\leq H'$) \unsure{supposons}}}

par \unsure{la|le} \unsure{propriété|propté} universelle de la passage
\diagram{\begin{tikzcd}[ampersand replacement=\&,column sep=0.55em,row sep=0.15em]
G_X \& \leq \& X\\
\rotatebox{90}{$\unsure{\preceq|\ll}$} \& \& \rotatebox{90}{$\preceq$}\\
\bar G_X \& \leq \& \unsure{F|P|\mathcal F}_X\\
\rotatebox{90}{$\leq$}\\
X
\end{tikzcd}}
On est ramené : prouver $X\leq G_X$

\deleted{D'où} \inserted{\unsure{Com|Comm.|Comme} \unsure{par|ça}} admettre
\deleted{$\unsure{X|G}\ \unsure{X|\ll}\ \unsure{\preceq|\leq}\ G$}
\margin{\deleted{$(\unsure{\leq|\ll|\preceq}F_X\unsure{\leq|<})$}}

\noindent\llap{\smash{\rule[-2.9\baselineskip]{.35pt}{3.4\baselineskip}}\hspace{.6em}}\deleted{\uline{C8} a) Soit \deleted{$X$ une multistrate, \deleted{\unsure{$X$|de|$G$}}. $\Phi$ \unsure{une}} $X$ une multistrate, $\Phi$ une subdivision de $X$, si $X\leq\Phi$ (i.e. $X$ est une multistrate de $\Phi$) on a $X=\Phi$}

\deleted{b) Si $X\preceq\Psi$, $X$ multistrate, $\Psi$ figure, alors $X=\Psi$}

\deleted{\unsure{Donc|Dém|Dem|Dans} dans le cas d'espèce on trouve}
\[
\deleted{X\ \deleted{\unsure{\leq|\preceq}}\ \leq\bar G_X\leq F_X=X}
\qquad\text{\deleted{donc}}\quad\deleted{X=\bar G_X}\quad\text{\deleted{, d'où}}
\]
\[
\deleted{X\preceq G_X\leq X}
\]

\scanpage{83}\mspage{82}

\noindent\llap{\smash{\rule[-4.9\baselineskip]{.35pt}{5.6\baselineskip}\hspace{.20em}\rule[-4.9\baselineskip]{.35pt}{5.6\baselineskip}}\hspace{.6em}}\uline{C 8}\textsubscript{(0)} Considérons les situations
\margin{\uline{provisoire !}}
\[
\begin{aligned}
\text{(a)}\qquad &X\leq\Phi\preceq X & \Phi&\in\mathcal F,\ X\in M\\
\text{(b)}\qquad &X\preceq\Psi\leq X & \deleted{\unsure{F|\Phi|\mathcal F}}\Psi&\in\mathcal F,\ X\in M
\end{aligned}
\]
\deleted{\unsure{C|A|On}}Alors on a $X=\Phi\quad\text{et}\quad X=\Psi$

Achevons la dém. Appliquant (a) à
\[
X\leq F_X\preceq X
\]
on trouve $X=F_X$, d'où, comme
\[
X\leq\bar G_X\leq F_X=X\quad,
\]
$X=\bar G_X$, d'où
\[
\deleted{\unsure{X=|X\leq|G_X=|G_X}}\bar G_X=X\preceq G_X\leq X
\]
et appliquant (b), on trouve $\deleted{\unsure{G|X|\Phi}}X=G_X$ et par suite $X\leq G_X$, qed.

\uline{Remarque} Bien sûr, C8 est nécessaire si on veut que $\ll$ soit une relation d'ordre. Et c'est aussi suffisant\unsure{,:|;:|;}

\uline{Cor} (de C8) \inserted{via prop. p\deleted{\unsure{g|.}}~76)} Si $F\ll G$ et $G\ll F$, on a $F=G$

\uline{Dém} \deleted{Soit $X\in\widetilde F$ \unsure{i.e.|ie.} $X\in M$, $X\leq F$} On va prouver $\unsure{F|\bar F}\leq G$\ ; \deleted{\unsure{(on|Car}} on aura par symétrie $G\leq F$, donc $F=G$.
\deleted{Il faut prouver \deleted{\unsure{pour|que}} pour $X\in M$, $(X\unsure{\leq|<}F)\ \unsure{-|\Rightarrow|=} X\leq G$.}

\diagram{\begin{tikzcd}[ampersand replacement=\&,column sep=0.45em,row sep=0.1em]
\deleted{F} \& \deleted{\ll} \& \deleted{G} \& \deleted{\ll} \& \deleted{F}\\
\rotatebox{90}{$\deleted{\leq}$} \& \& \rotatebox{90}{$\deleted{\leq}$} \& \& \rotatebox{90}{$\deleted{\leq}$}\\
\deleted{F_X} \& \deleted{\ll} \& \deleted{G_X} \& \deleted{\ll} \& \deleted{X}\\
\rotatebox{90}{$\deleted{\leq}$}\\
\deleted{X}
\end{tikzcd}}

\scanpage{84}\mspage{83}

On a par définition de $\ll$
\diagram{\begin{tikzcd}[ampersand replacement=\&,column sep=0.65em,row sep=0.15em]
F \& \leq \& G' \& \preceq \& G\\
\& \& \& \& \rotatebox{90}{$\geq$}\\
\& \& \& \& F'\\
\& \& \& \& \rotatebox{90}{$\succeq$}\\
\& \& \& \& \deleted{G}F
\end{tikzcd}}

on complète par C\deleted{\unsure{6|8}}7
\diagram{\begin{tikzcd}[ampersand replacement=\&,column sep=0.6em,row sep=0.15em]
F \& \leq \& G' \& \preceq \& G\\
\& \& \rotatebox{90}{$\deleted{\unsure{\leq|\preceq}}\geq$} \& \& \rotatebox{90}{$\geq$}\\
\& \& H \& \preceq \& F'\\
\& \& \& \& \rotatebox{90}{$\succeq$}\\
\& \& \& \& \deleted{G}F
\end{tikzcd}}

\deleted{\unsure{d|S}} donc
\begingroup
\makeatletter
\tagsleft@true
\let\veqno\@@leqno
\makeatother
\[
F\leq H\preceq F \tag*{$(*)$}
\]
On va en conclure $F=H$ (généralise C8 (a)) d.'où $G'=F$, $F'=F$, d'où
\margin{Mais ce sont \uuline{les \={m}s} cas, en échangeant les \unsure{rôles|rôle} de $F$ et $H$ on trouve $H\preceq F\leq H$ !}
\[
F\preceq G\leq F \tag*{$(**)$}
\]
\endgroup
on va en conclure $F=G$ (généralise C8 (b)).

Prouvons d'abord ce dernier pt : il \deleted{\unsure{faut}} suffit de prouver $F\preceq G$, \unsure{i.e.|ie.} \deleted{\unsure{Si|soit|Q|G}} $X\in M$, $X\leq F\Longrightarrow X\deleted{\unsure{<|\leq|\in|\preceq}}\leq G$, c'est là. \unsure{ess.|cor.|ex.|un.|m.} le corollaire p.~82.

\deleted{Pour \unsure{traiter|l'autre} $(\unsure{*|**})$}

On a essayé ici de se débrouiller aussi loin que possible, sans \unsure{critères|critère} et raisonnements de disjonction, genre C5.

\scanpage{85}\mspage{84}

\uline{Remarque} Soi\deleted{\unsure{en|n}}t $\mathcal F$ \deleted{\unsure{un|ens.|ens|une}} muni de deux \unsure{relations|relation} d'ordre $\leq$, $\preceq$, satisfaisant les conditions suivantes

a) Pour
\diagram{\begin{tikzcd}[ampersand replacement=\&,column sep=0.55em,row sep=0.25em]
F \& \deleted{\unsure{\leq|\preceq}}\geq \& G\\
\rotatebox{90}{$\preceq$}\\
F'
\end{tikzcd}}
on peut compléter de façon unique en
\diagram{\begin{tikzcd}[ampersand replacement=\&,column sep=0.65em,row sep=0.2em]
\& F \& \geq \& G\\
(*) \& \rotatebox{90}{$\preceq$} \& \& \rotatebox{90}{$\preceq$}\\
\& \deleted{\unsure{G|F}}F' \& \geq \& G'
\end{tikzcd}}

b) Pour
\diagram{\begin{tikzcd}[ampersand replacement=\&,column sep=0.55em,row sep=0.2em]
\deleted{F}G \& \leq \& F\\
\rotatebox{90}{$\preceq$}\\
G'
\end{tikzcd}}
, on peut compléter (de façon pas \unsure{néc\textsuperscript{t}|néc'.} unique) en
\diagram{\begin{tikzcd}[ampersand replacement=\&,column sep=0.65em,row sep=0.2em]
G \& \leq \& F\\
\rotatebox{90}{$\preceq$} \& \& \rotatebox{90}{$\preceq$}\\
G' \& \leq \& F'
\end{tikzcd}}

c) Si
\[
\begin{aligned}
&F\leq G\preceq F &&\text{(ou ce\unsure{.} qui revient au }\unsure{=|\bar m}\\
&G\preceq F\leq\deleted{\unsure{F|I}}G &&),
\end{aligned}
\]
alors $F=G$.

\uline{Alors} :
\deleted{(i) Posant $F\ll\ \overset{\mathrm{def}}{\Longleftrightarrow}\ \exists$}
(i) Posant $\deleted{F}G\ll\deleted{G}F\ \overset{\mathrm{def}}{\Longleftrightarrow}\ \exists\deleted{G}F'$ \unsure{avec|avc} $G\leq F'\preceq F$
la relation $\deleted{\unsure{\leq|<}}\ll$ est une relation d'ordre (et de préordre, si on ne suppose que a) b) )

\scanpage{86}\mspage{85}

(ii) Dans le cas a), le carré $(*)$ est cartésien pour $\ll$, i.e.
\[
G'=\underset{\ll}{\mathrm{Inf}}(F',G)
\]
\margin{Il manque \deleted{\unsure{peut-être}} \inserted{sans doute} des hypothèses ?? P.ex. que les $\mathrm{Inf}$ finis existent : pour $\leq$, et\ldots\ Il \uline{faudrait tirer} au \uline{clair la situation}}

\deleted{\unsure{(iii)|(ii)|(iv)}} (iii) Si $\deleted{\unsure{G|G'}}$
\diagram{\begin{tikzcd}[ampersand replacement=\&,column sep=0.7em,row sep=0.2em]
\deleted{\unsure{X|G|F}}F \& \deleted{\unsure{\leq|\preceq}}\geq \& G\\
\rotatebox{90}{$\leq$} \& \& \rotatebox{90}{$\deleted{\unsure{\preceq|\leq}}\leq$}\\
F' \& \geq \& G'
\end{tikzcd}}
est cartésien pour \rotatebox{90}{$\leq$}, il est cartésien pour $\deleted{\unsure{\leq|\preceq|\text{le}}}\ll$.

Dans les axiomes des \unsure{cartésies|cartésiens|cartiers|critères}, il faudrait séparer

\tikz[baseline=(n.base)]\node[draw,circle,inner sep=0.7pt](n){$\unsure{\alpha|*|K}$}; un paquet concernant $\leq$ tout seul
\inserted{(C1 \unsure{;|:|à} C4)}

\tikz[baseline=(n.base)]\node[draw,circle,inner sep=0.7pt](n){$\unsure{\beta|8|*|\varphi}$}; un paquet autour des propriétés ``formelles'' reliant $\leq$, $\preceq$ (\unsure{ess\textsuperscript{t}.|ess!|ord.|or!} celles qu'on vient d'esquisser) sans questions de support ni de disjonction, et en utilisant le minimum sur $\leq$ tout seul.

Conditions qui assurent, quand les sup majorés existent pour $\leq$, que ceux ci commutent à l'induction sur une ``subdivision.

Au besoin, si on \unsure{ne|en} veut se donner la peine de \deleted{\unsure{d|l}}dégager les axiomes minimum, donner le paquet des propriétés formelles ``évidentes'' du pt de vue signification géom,

\scanpage{87}\mspage{86}

\tikz[baseline=(n.base)]\node[draw,circle,inner sep=0.7pt](n){$*$}; Axiome de \uline{recollement} des subdivisions

\tikz[baseline=(n.base)]\node[draw,circle,inner sep=0.7pt](n){$*$}; \deleted{\unsure{A|U|Il|M}}Un paquet d'axiomes de disjonction

\deleted{\unsure{C6 C8|C8}}\qquad -\quad -\quad -

\uline{C8}\quad \uline{Axiome} \uline{de recollement} \uline{des subdivisions}

\noindent\llap{\smash{\rule[-2.5\baselineskip]{.35pt}{3.1\baselineskip}}\hspace{.6em}}Soit $\unsure{\mathbb F|F}= \operatorname{Sup} F_i$, et pour tt $i$, soit $F'_i$ une subdivision de $F_i$, telle que pour $i,j$, on ait $F'_i\mid F_i\cap F_j=F'_j\mid F_i\cap F_j$, alors $F'=\operatorname{Sup}F'_i$ existe et est une subdivision de $F$.

\deleted{Soit} (Devrait suivre C6\unsure{\}|)}, \unsure{Cela|cela} remplace C6 (c)

\uline{Cor} \unsure{Soient|Soit} \deleted{$F$ et $\unsure{F_a|F_i|F_\alpha}$, alors} \inserted{$F=\coprod F_i$, alors}
\[
\operatorname{Subdiv}(F)\simeq\prod_i\operatorname{Subdiv}(F_i)
\]

\uline{Axiomes de disjonction et de support}.

Pour trouver les axiomes sous forme aussi forte et sensible que possible, c'est le moment de \unsure{reintroduire|introduire|réintroduire} les lieux.

\deleted{\deleted{\uline{Definition} On appelle} \inserted{Prop.} Soit $F\in\mathcal F$, \inserted{$F\ne\phi_F$} Conditions equivalentes}

\deleted{a) $F'\preceq F\Longrightarrow F'=F$ \qquad i.e. $\deleted{\unsure{F|E}}F$ est minimal dans $\mathcal F$ pour $\preceq$}

\deleted{b) $F'\ll F$, \inserted{$F'\ne\phi$} $\Longrightarrow F'=F$ \qquad i.e. $F$ est minimal dans $\mathcal F$ \inserted{$\unsure{-|\setminus}\{\phi\}$} pour $\ll$}

\deleted{et ceci implique que $F\in M$\unsure{,} \deleted{\illegible{1}} et $F$ est minimal dans $\deleted{\unsure{\mathcal F|F}}M$ (\unsure{ou en|ou dans} $\deleted{\unsure{\mathcal F|F}}\unsure{M|\mathcal F}$) pour $\leq$}

\scanpage{88}
\mspage{87}
\deleted{\unsure{87|8|86}}

\deleted{Dém. $b)\Rightarrow a)$ trivial, prouvons $a\Rightarrow b$. Soit $F'\ll F$ i.e. $F'\leq\bar F\preceq F$. Par $\unsure{a|\text{--}}$, on a $\bar F=F$, donc $F'\leq F$. \deleted{\unsure{Par} $C7$, $\exists$} On a donc $F'\leq F$ \deleted{et $F'\unsure{\preceq|\leq}F$, \unsure{on veut|on} que ça implique $F'=F$.} \deleted{Mais si \unsure{le|la|l}} On va le prouver tantôt.}

\deleted{Ceci dit, comme $\ll$ \inserted{est} impliqué \inserted{par} $\leq$, il en resulte que (si $F\neq\phi$) $F$ est \deleted{minimal dans $\unsure{\mathcal F|\mathrm{Supp}}\unsure{\setminus\{\phi\}|\{\phi\}|\{x\}|\psi\{\phi\}}$ pour $\leq$ : Si $X'<$}}

\noindent\llap{\smash{\rule[-2.1\baselineskip]{.35pt}{2.8\baselineskip}}\hspace{.6em}}\uline{Définition} On appelle \uline{lieu} \inserted{de $\mathcal E$} un élément de $\mathcal F\setminus\{\phi_{\unsure{F|\mathcal F}}\}$ minimal pour $\ll$, i.e. une figure $\deleted{\unsure{X|x}}F$ telle que $\deleted{\unsure{X|x}}F\neq\phi_{\mathcal F}$, et $F'\ll F$, $F'\neq\phi_F$ implique $F'=F$.

Ainsi, on voit qu'un lieu est néc. $\in M$, et c'est aussi un élément de $M$ minimal pour $\leq$ (mais il n'est pas vrai, dans l'état actuel des choses, -- i.e. \unsure{avec|en un|en une} les axiomes qu'on \unsure{a} introduits jusqu'à présent) l'inverse soit vrai. (Ce sera vrai cependant pour les ``\unsure{contrées|cartésies|cartésiens|continu|continuum} en boules''\ldots)

On désigne par $\mathcal L$ l'ensemble des lieux.

\deleted{\unsure{On|De}} Pour toute $F\in\mathcal F$, on appelle \uline{ombre de $F$} l'ens. des $x\in\mathcal L$ tels que $x\ll F$, notation $\lvert F\rvert$, et \uline{multiombre} l'ens

\scanpage{89}
\mspage{88}
des parties $\deleted{\unsure{\widehat F|\widehat H}}$ de $\mathcal L$
\[
\widehat F=\bigl\{\,\lvert X\rvert\mid X\in\bar F\quad\text{i.e. }X\in M,\,X\leq F\,\bigr\}
\]

Ainsi on a
\[
\lvert\overset{\deleted{\unsure{\wedge|{-}}}}{F}\rvert\subset\mathcal L,\qquad
\widehat F\subset\mathcal P(L)
\]

et on a
\[
\lvert F\rvert=\deleted{\unsure{\mathcal L|\coprod}}\bigcup_{X\in\bar F}\lvert X\rvert=\bigcup_{A\in\unsure{\widehat F|\widehat{\bar F}}}A\,.
\]

Cette relation signifie que si $x\in\mathcal L$ est tel que $x\unsure{\deleted{\leq}\ll|\ll}F$, $\exists\,X\in\bar F$ tel que $x\ll X$, \unsure{\S|(*)|\&}
\margin{Mais pour \unsure{le|la} prouver, on a utilisé axiome d'invention}
et pour ceci, voir prop. page \unsure{78|76}. Il y a \unsure{même|\={m}} un plus petit tel $X$ (\unsure{=|:} Cor à la prop).

\uline{Proposition} La multi-ombre \inserted{$\Phi=$} $\widehat F$ d'une figure est une famille \uline{admissible} de parties de $\mathcal L$, i.e.
\[
\begin{aligned}
a)\quad &\forall\,\deleted{\unsure{a\in\mathcal L|a\in|x\in|A\in}}A\in\Phi,\ \text{posant }\ \unsure{\delta|\partial}A=\bigcup_{\substack{B\subset A\\B\neq A}}B\\
&\text{on a}\qquad \delta A\neq A\quad\text{\unsure{i.e|ie|ie.}}\quad A^{\unsure{\circ|o}}\overset{\mathrm{def}}{=}A\setminus\delta A\neq\phi
\end{aligned}
\]

\[
b)\quad \forall\,A,B\in\Phi,\ \text{on a}\qquad A\cap B=\bigcup_{\substack{C\in\Phi\\C\subset A\cap B}}C
\]

\cotefile{156-6}
\scanpage{2}\mspage{1}
\margin{\begin{tabular}{|l}GF \uline{VI}\\\hline\end{tabular}}
\begin{center}Analysis situs (\unsure{deuxième|deuxiême} mouture)\end{center}

\margin{18 juin 86}Avant de décrire ce qu'est une \uline{contrée}, je vais décrire ce qui sera une ``contrée avec notion de multistrates'' -- la famille des multistrates \inserted{choisies} jouant un peu le rôle \unsure{de'|d'|des} une famille d'ouverts \deleted{\unsure{basiques|basique} \unsure{pour}} \inserted{engendrant} une topologie, ou une famille génératrice d'éléments d'un topos. Je vais donc commencer par

\textcircled{I}\quad ``Algèbre de figures'' ou ``\uline{\deleted{\unsure{a}}Atelier}''.

1.\quad Une \uline{algèbre de figures} implique avant tout trois types d'objets, les \uline{lieux}, les \uline{multistrates}, les \uline{figures}, formant trois ensembles
\begingroup
\makeatletter
\tagsleft@true
\let\veqno\@@leqno
\makeatother
\[
\mathcal{L},M,\mathcal{F},\tag*{(1.1.)}
\]
\endgroup
liés entre eux par diverses applications, et chacun muni de diverses structures. Ainsi, on a des applications canoniques injectives
\diagram{\begin{tikzcd}[ampersand replacement=\&,column sep=large]
{\text{(1.2)}} \& \mathcal{L} \arrow[r,hook,"\underset{\unsure{\cdot}}{b)}"] \& M \arrow[r,hook,"a)"] \& \mathcal{F}{,}
\end{tikzcd}}
que nous utiliserons souvent pour identifier un lieu à une multistrate particulière, et une multistrate à une figure particulière -- $\mathcal{L}$ à un sous-ens. de $M$, $M$ \unsure{à|:} un sous-ens de $\mathcal{F}$.

Il y a d'autre part deu\deleted{\unsure{x}}x. autres paires d'applications, que voici.

\scanpage{3}\mspage{2}
\margin{b) se déduit de a) et de 1.2. a)}
\diagram{\begin{tikzcd}[ampersand replacement=\&,column sep=small,row sep=small,
execute at end picture={\draw (aa.north west) .. controls +(0,-.25) and +(.04,.16) .. (aa.south west) .. controls +(-.04,-.10) and +(-.16,.03) .. (bb.north west) .. controls +(.08,-.10) and +(-.02,.25) .. (bb.south west);}]
{\text{(1.3)}} \& |[alias=aa]| {\text{a)}} \& {\mathcal F\hookrightarrow\mathcal P(M)} \& {F\mapsto\widetilde F=\operatorname{\text{Dépl}}(F)\subset M} \\
\& |[alias=bb]| {\text{b)}} \& {\mathcal F\hookrightarrow\mathcal P(\mathcal P(M))} \& {\text{plus précisément }\mathcal F\to\operatorname{Fig}(M)} \\
\& {\substack{\text{ne sert}\\\text{à rien}}}\arrow[u,bend left=24]\arrow[rr,bend right=13] \& \& {F\mapsto\overset{\deleted{\unsure{\sim}}}{\mathrm F}\mathrm{ig}_{M}(F)\subset\mathcal P(M)}
\end{tikzcd}}
où $\operatorname{Fig}(M)$ désigne la partie de $\mathcal P(\mathcal P(\unsure{\cdot}M))$ formée des figures ensemblistes dans $M$. \deleted{On}

On peut considérer que la première application correspond à une relation entre $M$ et $\mathcal F$, appelée relation d'incidence. Pour une figure $F$, $\widetilde F$ s'appelle l'ensemble des \uline{multistrates} \uline{incidentes}, ou le \uline{déployement} de la figure $F$. Si $X\in M$, $F\in\mathcal F$, on dira que la multistrate $X$ est \uline{incidente} à la figure $F$, ou encore que c'est une \uline{strate de} \deleted{$F$}\inserted{\uline{la figure} $F$}, si $X\in\widetilde F$. D'autre part, tout\deleted{e} élément \inserted{$X$} de $M$ (ie. \deleted{toute} multistrate), regardé comme une figure par (1.2), admet un déployement $\widetilde X$, et \deleted{la multistrate} \deleted{\unsure{mais}} \deleted{\unsure{on|ou}} on pose
\margin{$\unsure{\bullet|\cdot}$}
\begingroup\makeatletter\tagsleft@true\let\veqno\@@leqno\makeatother
\[
\begin{aligned}
\text{\unsure{``}}\operatorname{Fig}_{M}(F)&=\{\widetilde X\mid X\in\widetilde F=\operatorname{Depl}(F)\}\unsure{\subset|\in}\mathcal P(M)\\
\text{ie.}\quad \operatorname{Fig}_{M}(F)&\in\mathcal P(\mathcal P(M))
\end{aligned}\tag*{(1.4)}
\]
\endgroup
et il résultera des axiomes que c'est une figure ensembliste dans $M$, indexée fidèlement par l'\unsure{ens|un} $\widetilde F$ des strates de $F$.

\deleted{\unsure{En fait}} \deleted{$\unsure{\mathcal F|F}$} En fait, $M$ sera muni d'une relation d'ordre $\leq$, cf plus bas, et $\widetilde F\inserted{\subset M}$ sera une partie fermée de \deleted{\unsure{la|de}} $M$, et pour tout $X\in\widetilde F$, \deleted{\unsure{$X$}} on aura
\begingroup\makeatletter\tagsleft@true\let\veqno\@@leqno\makeatother
\[
\widetilde X=\{Y\in M\mid Y\leq\unsure{x|X}\}\tag*{(1.5.)}
\]
\endgroup

\scanpage{4}\mspage{3}
À cause de cette interprétation, le passage de $\widetilde F\subset M$ à $\operatorname{Fig}_{M}(F)$ est à tel point évident, que \inserted{la \unsure{considération} de} cette figure ensembliste dans $M$ me semble \unsure{rarement|vraiment} importante -- mais à voir.....

\deleted{Quand on identi}

On pourra utiliser les deux applications (1.3) pour identifier une figure, suivant les cas, \unsure{à} une partie de $M$ [via $F\mapstochar\relbar\widetilde F$] ou à une figure ensembliste dans $M$ [via $F\mapsto\operatorname{\text{Dépl}}(F)$]. Dans le premier cas, \unsure{comparant|composant} \unsure{avec|à un} $M\unsure{\lhook\joinrel\relbar|\hookrightarrow}\mathcal{F}$ (1.2)\deleted{\unsure{(}}, par la formule 1,5. Cette dernière inclusion s'identifie à
\[
X\mathrel{\mapstochar\relbar\mapstochar}\widetilde X=\{Y\in M\mid Y\leq X\}
\qquad M\to\mathcal{P}_{\mathrm{ferm}}(M).
\]

Donc il faudra être prudent alors \inserted{\unsure{avec}} l'identific\unsure{ation} dont il a été question, des\deleted{\unsure{une|un}} $\unsure{X|X'}\in M$ \unsure{ou|en} \unsure{multistrates|multistrate} à des figures particulières, quand celles-ci sont interprétées comme des parties de $M$: La figure associée à $X\in M$ n'est nullement $\{X\}$, mais $\widetilde X\rbrace$ \deleted{\unsure{$f$|$\mathcal F$|$F$|l'}} ens des multistrates majorées par $x$ pour $\leq$ (ou encore, des sous-multistrates, ou multistrates incidentes à $X$, ou encore des \deleted{\unsure{$f$|$\mathcal F$|$F$}} \uline{strates} de $x$).

\scanpage{5}\mspage{4}
D'autre part, il y a deux \inserted{\uline{\unsure{autres}}} applications
\begingroup
\makeatletter\tagsleft@true\let\veqno\@@leqno\makeatother
\[
\left\{\begin{aligned}
\text{a)}\quad &\unsure{\mathcal{F}|\mathcal{E}}\longrightarrow\mathcal{P}(\deleted{\unsure{\mathcal{F}|F}}M)&
 F&\longmapsto\operatorname{Omb}_{\unsure{M|\deleted{\mathcal F}M|m}}(F)\\
\text{b)}\quad &\mathcal{F}\hookrightarrow\mathcal{P}\mathcal{P}(\deleted{\unsure{\mathcal{F}|F}}M)&
 \text{et }=\quad\unsure{\mathcal{F}|\widetilde{\mathcal{F}}}&\longrightarrow\operatorname{Fig}(M\unsure{]|)}\subset\mathcal{P}\mathcal{P}\unsure{M|(M)}\\
&& F&\mapsto\deleted{\unsure{O}}\operatorname{Mul\unsure{t|\deleted{\unsure{f}}t}omb}_{M}(F),
\end{aligned}\right.\tag*{(1.\deleted{\unsure{4}}6)}
\]
\endgroup
où $\operatorname{Omb}_{M}(F)$ s'appelle l'\uline{\unsure{ombre|ogombre}} \uline{de la} figure $F$ sur $M$, ou la \uline{grande} \uline{ombre} \uline{de} $\mathcal{F}$. \deleted{\unsure{La}}Une multistrate \inserted{$X$} qui est dans la grande ombre d'une figure $F$ est dite ``\uline{raffiner}'' \uline{la figure} $F$ (c'est cette fois déduit d'une relation d'ordre $\ll$ sur $\mathbb{F}$, le raffinement, par la condition $X\ll F$). Il n'est pas clair que cette application $\mathcal{F}\unsure{\relbar|\to}\mathcal{P}(M)$ soit injective, ie. que $\mathcal{F}$ s'identifie à un ens de parties de $M$, via cette application. Par contre il en sera ainsi de (1.6) b), qui est déduit de a) via \deleted{\unsure{$F$|$\mathcal F$}}la \unsure{formation|fonction} de $\widetilde F$, comme (1.3 b) est déduit de (1.3. a):
\begingroup
\makeatletter\tagsleft@true\let\veqno\@@leqno\makeatother
\[
\operatorname{Multomb}(F)=\{X\in\widetilde F\mid\operatorname{Omb}X\},\tag*{(1.\deleted{\unsure{5|6}}7)\quad\deleted{(4.5)}}
\]
\endgroup
alors que $\operatorname{Omb}(X)$ s'explicite par
\begingroup
\makeatletter\tagsleft@true\let\veqno\@@leqno\makeatother
\[
\unsure{\operatorname{Multomb}|\operatorname{MultOmb}}(X)=\{Y\in\deleted{\unsure{\mathcal{F}|F}}M\mid Y\ll X\},\tag*{(1.\deleted{\unsure{6|7}}8)}
\]
\endgroup
analogue à (1.5), où cette fois on a utilisé la relation \inserted{d'ordre} $\ll$ ald $\leq$ -- \deleted{C'est \unsure{la|le}}
